\scanpage{29}
\mspage{15}
\textcircled{9} \uline{Une conjecture arithmético-géométrique}

Soient $g,\nu\geq0$ des entiers. On désigne par $M_{g,\nu}$ le topos modulaire \inserted{sur $\mathbf Q$} pour les courbes \inserted{relatives} lisses propres relativement connexes de genre $g$, \uline{munies} de \deleted{\unsure{un}} $\nu$-uple de \unsure{$\nu$ sections|sections} disjointes sur chaque fibre. Sauf pour $g=0$, $\nu=0,1,2$ et $g=1$, $\nu=0$, ce topos est associé à une étendue algébrique, \unsure{notion proche = schémable}.

On a un hom naturel \deleted{\unsure{extér.}}
\[
\pi_1(M_{g,\nu};\xi)\longrightarrow\operatorname{Autextloc}^{\circ}\bigl(\pi_1(\overline X\setminus\{x_1,\ldots,x_\nu\})\bigr)\tag{1}
\]
où $(X;x_1,\ldots,x_\nu)$ est la courbe \inserted{(sur $\overline{\mathbf Q}$)} de genre $g$ avec $\nu$ pts marqués associée à $\xi$ \inserted{$\in M_{g,\nu}(\overline{\mathbf Q})$}. \deleted{\unsure{On}} L'exposant $\circ$ marque le fait qu'on prend les autom de $\pi_1$ qui n'échangent pas entre eux les classes de Cartan associés aux $x_i$.
\margin{il faut peut-être être plus \unsure{strict}, et exiger la commutation \unsure{:} des opérations \unsure{:} \unsure{définies} de $\mathfrak S_\nu$ sur le $\pi_1$ ?? Il faut regarder surtout le cas $g=0,\nu=3$\ldots}

Notons qu'on a aussi
\[
\pi_1(\overline M_{g,\nu};\unsure{\xi|j})\longrightarrow\operatorname{Autextloc}^{\circ\circ}\bigl(\pi_1(X\setminus\{x_1,\ldots,x_\nu\})\bigr)\tag{2}
\]
où le double $\circ$ indique qu'on prend des autom \inserted{\uline{stricts}} compatibles avec la structure \unsure{locale}, i.e. induisant l'identité sur $T=\operatorname{Tate}(\overline{\mathbf Q})$\ldots
\margin{Aut. \uline{\unsure{orientés}}}
\[
\pi_1(\unsure{M_{g,\nu}^{an}|\overline M_{g,\nu}};\unsure{\xi|j})\longrightarrow
\underbrace{\operatorname{Autextloc}^{\circ\circ}\Pi_1(\unsure{X^{an}|\overline X}\setminus\{x_1,\ldots,x_\nu\})}_{=\Gamma_{g,\nu}\ (\text{Teichmüller})}\tag{2 bis}
\]
dont on sait (théorie top.) que c'est un
\mspage{28}
